% Borrador para IEEEtran. Compilar con pdflatex dos veces o abrir en Overleaf.
% No contiene los documentos institucionales privados.
\documentclass[conference]{IEEEtran}
\usepackage[utf8]{inputenc}
\usepackage[T1]{fontenc}
\usepackage[spanish]{babel}
\usepackage{amsmath}
\usepackage{url}
\usepackage{hyperref}
\title{Yachi: diseño de un asistente institucional trazable para procedimientos universitarios}
\author{\IEEEauthorblockN{Autores del equipo: por completar}
\IEEEauthorblockA{Universidad Yachay Tech\\Urcuquí, Ecuador}}
\begin{document}
\maketitle
\begin{abstract}
Se presenta el diseño e implementación de Yachi, un prototipo conversacional orientado a procedimientos, representación estudiantil y ubicaciones universitarias. Combina recuperación de información institucional, fuentes visibles, guías estructuradas y navegación del campus. Se distinguen el simulador de interfaz, el modo extractivo y la generación asistida por modelo. La verificación funcional incluye compilación, contratos HTTP y recorridos de navegador con API aislada. La validación factual bilingüe y el estudio con usuarios permanecen pendientes; no se reportan porcentajes de precisión ni impacto sin una medición correspondiente. Este documento es un borrador de trabajo, no una investigación concluida.
\end{abstract}
\begin{IEEEkeywords}
asistente institucional, recuperación de información, trazabilidad, experiencia de usuario, evaluación
\end{IEEEkeywords}
\section{Introducción y problema}
La dispersión de reglamentos, calendarios y procedimientos dificulta que un estudiante identifique qué gestión realizar, en qué orden, con qué documentos y ante qué dependencia. El Reto 1 del ideatón solicita un prototipo o propuesta con información verificable, español e inglés y una experiencia intuitiva. Se delimita el problema a orientación informativa: no se realizan matrícula, pagos, modificación de notas ni firma electrónica.

La pregunta de investigación propuesta es: ¿en qué medida una conversación con fuentes verificables y guía estructurada mejora la resolución correcta de consultas institucionales frente a buscar directamente en documentos? La hipótesis debe contrastarse con usuarios; no se considera demostrada por la existencia del prototipo.
\section{Estado del arte}
Pounce documenta atención por mensajería y evaluación experimental en procesos de matrícula\cite{pounce}. Genie describe acceso a servicios personales y navegación\cite{genie}. Sunny documenta orientación bilingüe para públicos específicos\cite{sunny}; LOLA se orienta a preinscripción y matrícula\cite{lola}. Estas publicaciones institucionales acreditan funciones descritas en sus fechas, no una auditoría actual de todos sus servicios.

La recuperación aumentada con generación conecta el modelo con información recuperada\cite{rag}. En un entorno institucional, recuperar y citar un documento no prueba por sí solo que la respuesta sea correcta: deben evaluarse correspondencia factual, vigencia, cobertura y rechazo adecuado. La aportación propuesta de Yachi es integrar evidencia visible, guía de procedimiento y ubicación; no se reivindica exclusividad de sus componentes.
\section{Metodología}
Se propone un ciclo de inventario documental, estructuración de metadatos, implementación, evaluación y corrección. El corpus oficial se administra por un canal autorizado y se excluye del repositorio público. Cada documento requiere fuente, tema, versión, fecha, acceso y vigencia. El equipo debe completar el inventario y la revisión humana, incluyendo representatividad estudiantil.

La evaluación distinguirá tres configuraciones: simulador, backend extractivo y backend con generación. Las métricas de una no se atribuirán a otra. Las pruebas con respuestas HTTP simuladas verifican estados de interfaz y recuperación de errores, no calidad del modelo ni persistencia de producción.
\section{Diseño e implementación}
La interfaz utiliza React y TypeScript; la API, FastAPI y LangGraph, con PostgreSQL y recuperación textual/vectorial. El enrutamiento incluye consultas ambiguas, ausencia de información, restricciones de ámbito y localización. Las fuentes visibles proceden de metadatos del servidor. Las guías curadas dependen de evidencia del corpus y requieren revisión institucional.

El chat presenta estados breves durante el procesamiento y texto progresivo. Al finalizar se elimina el historial técnico de actividad y se conservan respuesta y fuentes. La presentación progresiva puede repartir contenido ya recibido; no debe confundirse con una medición de velocidad de inferencia. Se permite detener, se respeta movimiento reducido y no se debe desplazar a quien lee mensajes anteriores.

El portal incorpora documentos, temas, tendencias, consultas sin resolver, filtros, validación, costos y consultas registradas. Estas últimas cuentan con estados de atención y bitácora; el contacto humano es manual, sin envío automático de correos. Los adjuntos del público se deshabilitan hasta disponer de análisis real, manteniendo la carga administrativa de documentos.
\section{Validación y resultados disponibles}
En esta fase pasan seis pruebas aisladas de contratos HTTP, la compilación TypeScript/Vite y una regresión Playwright de menú móvil, foco, errores de feedback/evaluación, importación CSV, formulario móvil y atención de consultas. La compilación advierte sobre el tamaño de un paquete JavaScript. Estos resultados demuestran únicamente las condiciones cubiertas por cada prueba.

Se propone un conjunto interno de al menos 100 preguntas balanceadas ES/EN, con respuestas esperadas revisadas por personas, consultas ambiguas, fuera de ámbito y sin información. La cantidad no proviene de una obligación numérica del reto. Las preguntas deben separarse de las usadas para ajustar el sistema; los casos de un mismo procedimiento no deben producir una falsa independencia estadística.

Una comprobación adicional sobre una tabla temporal de PostgreSQL verificó el guardado del set y la reversión completa de un fallo intermedio. Esta prueba no modificó el set persistente de la aplicación ni sustituye la validación factual.

Para cada pregunta se conservarán configuración, fecha, pregunta, respuesta esperada, respuesta obtenida, documento, resultado y latencia. Se reportarán exactitud factual con rúbrica humana, coincidencia de fuente, rechazo apropiado y latencias p50/p95. La coincidencia de documento y la alerta de texto sin cita son indicadores automáticos parciales. La tabla de resultados empíricos queda pendiente: no se sustituye por valores del simulador.

En una prueba de tareas con estudiantes se propone medir resolución correcta, tiempo, errores y percepción de claridad, comparando el asistente con búsqueda documental. El protocolo, participantes y resultados deberán completarse con consentimiento y datos minimizados.
\section{Costo computacional}
Sean $S$ estudiantes, $q$ consultas por estudiante y día, $h$ fracción de aciertos de caché y $c_m$ costo por consulta sin acierto. El costo mensual estimado es:
\begin{equation}
 C = C_{infra} + 30Sq(1-h)c_m + C_{escritura}.
\end{equation}
El costo $c_m$ incluye entrada, salida y clasificación según la arquitectura elegida. La calculadora utiliza precios documentados de Anthropic\cite{pricing} y un servidor de referencia de DigitalOcean\cite{infra}. El supuesto de 4 vCPU y 8 GiB a USD 48 mensuales no constituye una medición de capacidad. La caché debe verificarse en la configuración real; no se presentarán ahorros supuestos como consumo observado.

Para $N$ vectores de dimensión $d$ con float32, el almacenamiento vectorial mínimo es $4Nd$ bytes, al que se añaden texto, metadatos e índices. El esquema actual usa dimensión 384. CPU, RAM, disco, concurrencia, respaldos y costos operativos requieren medición antes de una adopción institucional.
\section{Uso de IA}
OpenAI Codex asistió en código, correcciones, auditoría con Playwright, investigación y borradores. El usuario definió identidad, preferencias visuales y prioridades. La declaración del repositorio debe completarse con herramientas históricas confirmadas, responsables humanos y revisión de contenido. La aprobación visual no acredita revisión de cada hecho ni autoría humana de todo el código.
\section{Conclusiones y limitaciones}
Se dispone de un prototipo con fuentes visibles, conversación, mapas y herramientas administrativas. La viabilidad propuesta depende de documentos vigentes, responsables de atención y evaluación factual ES/EN. No se certifican ausencia universal de alucinaciones, conformidad WCAG, generación bilingüe completa ni funcionamiento de voz con hardware real. Kichwa permanece como valor agregado pendiente de validación competente.

Antes de entregar: completar autores, inventario, revisión de referencias, resultados empíricos, responsable institucional, revisión del aviso y ensayo del pitch. Publicar código no equivale a desplegar un servicio de producción.
\begin{thebibliography}{9}
\bibitem{pounce} Georgia State University, ``Reduction of Summer Melt,'' [En línea]. Disponible: \url{https://success.gsu.edu/reduction-of-summer-melt/}. Consulta: 20 oct. 2026.
\bibitem{genie} Deakin University, ``Deakin's Genie: a virtual digital assistant out of the bottle,'' [En línea]. Disponible: \url{https://www.deakin.edu.au/about-deakin/news-and-media-releases/articles/deakins-genie-a-virtual-digital-assistant-out-of-the-bottle}. Consulta: 20 oct. 2026.
\bibitem{sunny} Arizona State University, ``ASU launches bilingual chatbot to reach parents who want to go to college,'' 2020. [En línea]. Disponible: \url{https://news.asu.edu/20200709-asu-launches-bilingual-chatbot-reach-parents-who-want-go-college}.
\bibitem{lola} Universidad de Murcia, ``LOLA Chatbot (Corto),'' 2018. [En línea]. Disponible: \url{https://clutv10.um.es/pro/video?id=116011}.
\bibitem{rag} P. Lewis et al., ``Retrieval-Augmented Generation for Knowledge-Intensive NLP Tasks,'' NeurIPS, 2020. [En línea]. Disponible: \url{https://arxiv.org/abs/2005.11401}.
\bibitem{pricing} Anthropic, ``Pricing,'' [En línea]. Disponible: \url{https://platform.claude.com/docs/en/about-claude/pricing}. Consulta: 20 oct. 2026.
\bibitem{infra} DigitalOcean, ``Droplet Pricing,'' [En línea]. Disponible: \url{https://www.digitalocean.com/pricing/droplets}. Consulta: 20 oct. 2026.
\end{thebibliography}
\end{document}
